\documentclass[10pt,a4paper]{amsart}
\usepackage[margin=19mm]{geometry}
\usepackage{amsmath,amssymb,mathtools}
\usepackage[scr=rsfs,cal=euler]{mathalfa}
\usepackage{xfrac}
\usepackage[unicode,hidelinks,bookmarks=false]{hyperref}
\newcommand{\ie}{i.e.\ }
\newcommand{\cf}{cf.\ }
\newcommand{\resp}{respectively\ }
\newcommand{\Subsection}{\subsection}
\providecommand{\nobreakdash}{\nobreak\hbox{-}}
\setlength{\emergencystretch}{2em}
\multlinegap0pt
\usepackage[shortlabels]{enumitem}
\usepackage{amssymb}
\usepackage{xcolor}
\usepackage{algorithm}\usepackage{algpseudocode}
\usepackage{tikz}
\newtheorem{corollary}{Corollary}
\newcommand{\Fcal}{\mathcal{F}}
\newcommand{\Lcal}{\mathcal{L}}
\newcommand{\R}{\mathbb{R}}
\newcommand{\dist}{\operatorname{dist}}
\title{First-Order Optimality Conditions for Mathematical Programming with Equilibrium Constraints}
\author{Louis Shuo Wang}
\date{Source: arXiv:2605.00388v1 (CC BY 4.0)}
\begin{document}
\maketitle

\noindent\emph{Source and licence.} First-Order Optimality Conditions for Mathematical Programming with Equilibrium Constraints. Version arXiv:2605.00388v1 (CC BY 4.0), distributed by arXiv under the Creative Commons Attribution 4.0 International licence, http://creativecommons.org/licenses/by/4.0/, as recorded in the arXiv licence metadata for this exact version and shown on the abstract page https://arxiv.org/abs/2605.00388v1. Full licence notice: \texttt{licenses/arxiv-2605.00388-CC-BY-4.0.txt}.\par\smallskip

\noindent\textit{Editorial scope.} Counted unit: the article's corollary cor:331 (source0.tex lines 776--801). The article's complete original proof of it is delivered with it (source0.tex lines 803--818, one \texttt{proof} environment of 16 lines). No mathematics is written, repaired or re-derived: the statement and the proof are the article's own lines, in the article's own notation, and no retained line is substituted or re-flowed.  No further source-local context is needed: every cross-reference of every retained line resolves inside the two delivered spans.  Nothing was omitted from any retained span, so the package carries no omission record.  No retained line contains a citation, so the wrapper carries no bibliography and no bibliography entry.  No retained line contains a figure environment, an image-inclusion command or drawing code, so no image is delivered and the package declares no asset dependency.  Editorial formatting only, in addition: an AMS \texttt{amsart} preamble that restates the article's own declarations and macros used by the retained lines, namely the environments \texttt{corollary} on separate counters, together with the standard packages those lines need.  The retained environments print in this document as Corollary 1 in order of appearance, whereas the article prints the same items with the numbering of its own full document.  The complete native source of the article as distributed in the arXiv e-print for this exact version is bundled unchanged as \texttt{sources/199557/source0.tex}.
\begin{corollary}[Equivalent primal characterizations of stationarity]
\label{cor:331}
Assume SBCQ holds at $\bar z\in\Fcal$ and the full CQ holds at $\bar z$,
that is,
\[
\Lcal(\bar z;\Fcal)=T(\bar z;\Fcal).
\]
Then the following are equivalent:
\begin{enumerate}[label=\textnormal{(\alph*)}]
\item $\bar z$ is a stationary point of the MPEC;
\item
\[
dz\in \Lcal(\bar z;\Fcal)\Longrightarrow \nabla f(\bar z)^Tdz\ge 0;
\]
\item $dz=0$ is an optimal solution of
\[
\min_{dz}\ \nabla f(\bar z)^Tdz
\qquad\text{s.t.}\qquad dz\in \Lcal(\bar z;\Fcal);
\]
\item for every scalar $\alpha\ge \|\nabla f(\bar z)\|$ and every
$dz\in\R^{n+m}$,
\[
\nabla f(\bar z)^Tdz+\alpha\,\dist(dz,\Lcal(\bar z;\Fcal))\ge 0.
\]
\end{enumerate}
\end{corollary}

\begin{proof}
The equivalence of (a) and (b) follows directly from the definition of
stationarity and the equality $T(\bar z;\Fcal)=\Lcal(\bar z;\Fcal)$. The
equivalence of (b) and (c) is immediate. The implication (d)$\Rightarrow$(c)
is obvious. For (c)$\Rightarrow$(d), choose $dz'\in\Lcal(\bar z;\Fcal)$ with
\[
\|dz-dz'\|=\dist(dz,\Lcal(\bar z;\Fcal)).
\]
Then
\[
\nabla f(\bar z)^Tdz+\alpha\,\dist(dz,\Lcal(\bar z;\Fcal))
\ge
\nabla f(\bar z)^Tdz'+(\alpha-\|\nabla f(\bar z)\|)\|dz-dz'\|\ge 0.
\qedhere
\]
\end{proof}

\end{document}
